\documentclass[11pt,a4paper]{article}
\usepackage[utf8]{inputenc}
\usepackage[T1]{fontenc}
\usepackage{mathpazo}
\usepackage[a4paper,margin=20mm]{geometry}
\usepackage[expansion=false]{microtype}
\usepackage[hidelinks]{hyperref}
\usepackage{needspace}
\pagestyle{empty}
\setlength{\parindent}{0pt}
\setlength{\parskip}{4pt}
\setlength{\emergencystretch}{2em}
\newcommand{\parte}[2]{\par\addvspace{7pt}\Needspace{5\baselineskip}\noindent\textbf{#1}\hfill{\small #2}\par\nobreak}
\hypersetup{pdftitle={Resumen del diálogo sobre investigación teórica en Panamá},pdfauthor={Resumen de la presentación de Danesh Morales y Marc Urcia}}

\begin{document}
{\large\bfseries Diálogo sobre el desarrollo de oportunidades\\ de investigación teórica en Panamá}\par
{\small Danesh Morales y Marc Urcia \enspace\textbullet\enspace 31 de octubre de 2026}

\textbf{Resumen.} La charla propone fortalecer la investigación teórica en física y matemática en Panamá mediante proyectos accesibles, mentoría y colaboración nacional e internacional. Se muestran formas de comenzar desde la licenciatura, ejemplos de tesis y oportunidades en información cuántica. La discusión posterior aborda la preparación matemática, la búsqueda de supervisores, los experimentos remotos y el uso de inteligencia artificial. La propuesta concreta de los ponentes es acompañar a estudiantes en proyectos que puedan convertirse en tesis y, cuando los resultados lo permitan, publicaciones.

{\small Los tiempos son aproximados y corresponden al audio proporcionado, de 1 h 13 min 40 s. La exposición ocupa los primeros 35 minutos; después comienza el diálogo.}

\parte{1. Contexto y propósito}{00:00--04:15}
Se presenta la investigación teórica como complemento de los grupos experimentales del país. Su menor dependencia de infraestructura material permite iniciar muchos proyectos con una computadora, bibliografía y acceso a internet. Se destaca la posibilidad de colaborar a distancia, formar especialistas e involucrar a estudiantes desde segundo o tercer año de licenciatura mediante mentoría de profesores o investigadores en formación.

\parte{2. Áreas de investigación y aplicaciones}{04:15--06:10}
Se ilustran cuatro áreas: teoría de control y sistemas dinámicos para redes de agua, energía y riego; geometría diferencial para planificar movimientos de robots; materia condensada para estudiar arquitecturas de qubits semiconductores; e información cuántica para corregir errores. Los ejemplos muestran cómo problemas matemáticos o físicos teóricos pueden conectarse con aplicaciones y necesidades del país.

\parte{3. Cómo comenzar y qué puede producir una tesis}{06:10--12:30}
Marc describe cuatro vías: revisar y sintetizar literatura, reproducir y extender simulaciones, estudiar casos particulares de problemas abiertos y analizar los límites de los modelos. Se presentan trabajos sobre flujos de Couette--Taylor, órbitas en Schwarzschild, estimación de fidelidad mediante \emph{Circuit Benchmarking} y formas diferenciales en física clásica. Una tesis puede aportar resultados numéricos, análisis matemático o una exposición pedagógica rigurosa; las simulaciones también pueden motivar futuros experimentos.

\parte{4. Información cuántica y colaboración en Panamá}{12:30--22:55}
La computación cuántica se presenta como un punto de encuentro entre física, matemática, química e ingeniería. Se propone conectar estudiantes y grupos de distintas universidades y regiones. El vínculo con C-TASC y los semiconductores se plantea como una posibilidad que requiere estudiar capacidades y formar especialistas. Marc explica la \textbf{caracterización de errores}, que busca conocer el ruido y evaluar la fiabilidad de los procesos; Danesh explica la \textbf{corrección de errores}, que busca proteger y recuperar la información codificada.

\parte{5. Formación, financiación y oportunidades}{22:55--35:30}
Se aconseja contactar posibles supervisores para conocer sus proyectos, plazas y financiación antes de solicitar admisión: una beca externa no siempre es indispensable. También se recomienda seguir artículos mediante arXiv, SciRate y Semantic Scholar, y relacionar la tesis con investigación reciente. Se mencionan PSI Start, Journeys into Theoretical Physics, USEQIP, Qiskit Global Summer School, Outsourcing Math Research y la Escuela de matemáticas cuánticas. Al cerrar la exposición, los ponentes ofrecen orientación y proponen un grupo de contacto para explorar tesis de licenciatura.

\parte{6. Qué aprender para comenzar}{35:30--45:05}
En las preguntas se explica que el álgebra lineal permite empezar a trabajar con qubits y sistemas de dimensión finita. Otras herramientas, como teoría de grupos, dependen de la especialidad y pueden aprenderse progresivamente. Se recomiendan los cursos de IBM y los libros de John Watrous y de Nielsen y Chuang. Para quienes prefieran otras áreas, se señalan ecuaciones diferenciales, sistemas dinámicos y control, con aplicaciones en epidemiología y robótica.

\parte{7. Proyectos y posible asesoría desde Panamá}{45:05--50:40}
Danesh y Marc expresan disposición a orientar proyectos y explorar una coasesoría coordinada con docentes locales. Proponen partir de un artículo y modificar una hipótesis, el modelo de ruido o el número de qubits, comparando resultados y buscando una explicación matemática. Se plantea complementar simulaciones con ejecuciones remotas en equipos de IBM. La formalización de la asesoría, el acceso al hardware y la financiación se presentan como asuntos que habría que concretar.

\parte{8. Cómo contactar investigadores}{50:40--57:50}
Se distinguen tres motivos para escribir: consultar una duda específica sobre un artículo, buscar una plaza de posgrado o explorar una pasantía. El mensaje debe presentar al estudiante, explicar su trabajo e interés y formular una petición concreta. Se recomienda perseverar ante la falta de respuesta y aprovechar reuniones virtuales y comunidades académicas. Trabajar sobre un artículo puede facilitar una conversación con sus autores y demostrar iniciativa.

\parte{9. IA como apoyo a la investigación}{57:50--61:45 y 63:50--69:30}
Los participantes discuten su experiencia usando IA para bibliografía, programación, cálculos y revisión de argumentos. El consejo central es conservar el criterio propio: decidir qué pregunta importa, qué se quiere comunicar y cómo verificar los resultados. Se insiste en revisar cada paso antes de continuar, porque los errores pueden acumularse. Se desaconseja delegar sin orientación la generación de un proyecto o de un artículo completo; la mentoría y el juicio del investigador siguen siendo fundamentales.

\parte{10. Continuidad y transición hacia un posgrado}{61:45--63:50 y 69:30--73:40}
Se propone una próxima reunión más específica sobre áreas de la computación cuántica, problemas abiertos y proyectos posibles. Marc explica que su tesis de licenciatura no trató de cuántica: se acercó al área mediante cursos y comunidades, y continuó su formación durante la maestría. Su experiencia muestra que una tesis en otro tema no impide explorar una nueva especialidad. El encuentro termina con la intención de compartir las diapositivas y mantener el contacto para preguntas y colaboraciones.

\par\addvspace{8pt}
{\small\textbf{Fuentes.} Audio del encuentro del 31 de octubre de 2026 y presentación de apoyo (28 diapositivas), proporcionados para este resumen. \href{https://www.youtube.com/watch?v=kJAX_0uai8Y}{Grabación en YouTube}.}
\end{document}
